% Generated draft: requires mathematical and source review.
\documentclass[11pt]{article}
\usepackage[margin=25mm]{geometry}
\usepackage{fontspec}
\setmainfont{Latin Modern Roman}
\usepackage{amsmath,amssymb,mathtools,unicode-math}
\setmathfont{Latin Modern Math}
\usepackage{enumitem,needspace,xurl,hyperref}
\setlength{\parindent}{0pt}
\setlength{\parskip}{5pt}
\setlength{\emergencystretch}{3em}
\newcommand{\N}{\mathbb{N}}
\newcommand{\Z}{\mathbb{Z}}
\newcommand{\Q}{\mathbb{Q}}
\newcommand{\R}{\mathbb{R}}
\newcommand{\C}{\mathbb{C}}
\newcommand{\F}{\mathbb{F}}
\newcommand{\A}{\mathbb{A}}
\newcommand{\PP}{\mathbb{P}}
\newcommand{\E}{\mathbb{E}}
\newcommand{\eps}{\varepsilon}
\newcommand{\dd}{\,\mathrm{d}}
\begin{document}
\section*{UnsolvedMath ID: 800108}
\textbf{Source problem code:} catalogue entry
\subsection{Definitions and mathematical statement}
Identify the circle \(S^1\) with the ideal boundary of the hyperbolic disk \(\mathbb H^2\). A geodesic lamination is a closed subset of the space of unordered pairs of distinct points of \(S^1\), such that the corresponding hyperbolic geodesics have pairwise disjoint interiors. Its leaves are these geodesics, and its gaps are the connected components of their complement in \(\mathbb H^2\). The lamination is \emph{very full} if every gap is a finite-sided ideal polygon, meaning that its vertices lie on \(S^1\) and its sides are hyperbolic geodesics.

A circle action of a group \(G\) is a homomorphism \(\rho:G\to\operatorname{Homeo}(S^1)\). A lamination \(\Lambda\) is invariant under this action if, for every \(g\in G\) and every leaf with endpoints \(\{x,y\}\), the pair \(\{\rho(g)(x),\rho(g)(y)\}\) is also the endpoint pair of a leaf of \(\Lambda\). The action is faithful if \(\rho\) is injective. The group \(G\) is \emph{residually finite} if
\[
\forall g\in G\setminus\{1_G\}\quad
\exists\text{ a finite group }F\quad
\exists\text{ a homomorphism }q:G\to F
\quad q(g)\ne 1_F.
\]

\textbf{Source request, with unresolved action conventions.} The supplied record attributes to Calegari the conjectural assertion that a laminar group preserving a very full lamination is residually finite. Its concluding request asks that a group acting on the circle with a very full invariant lamination be shown to satisfy the finite-separation condition above.

The record does not define ``laminar group'' or specify whether the circle action must be faithful or orientation preserving. If faithfulness is intended, the action hypothesis must include injectivity of \(\rho\); if orientation preservation is intended, its image must lie in \(\operatorname{Homeo}^{+}(S^1)\). Neither convention is silently imposed here. Consequently, the admissible class of actions remains unresolved, and the source request is not asserted as a theorem for arbitrary, possibly nonfaithful actions. The record supplies no countability or finite-generation hypothesis.
\subsection{Short English statement}
The supplied question, attributed to Calegari, asks whether preserving a circle lamination whose gaps are all finite-sided ideal polygons forces the acting group to be residually finite. The source leaves faithfulness and orientation conventions unspecified, so these hypotheses require clarification before the intended conjecture has a definitive formulation.
\subsection{Sources}
\begin{itemize}
\item Local UnsolvedMath registry, supplied complete source record: permanent ID \(800108\), source problem code \texttt{catalogue entry}, title ``Foliations and Diffeomorphism Groups --- Question 49 (Calegari).'' The registry labels the problem \emph{open}; this is supplied status metadata, not an independently verified mathematical assessment.
\item Nariman and Yazdi, \emph{Problem List on Foliations and Diffeomorphism Groups} (2025), Question 49. Supplied source link: \url{https://arxiv.org/abs/2503.18274}. The record reports access on 2026-10-16 and extraction by \texttt{arxiv\_latex}; these provenance details are supplied and have not been independently verified here.
\item Attribution and context in the supplied record: Calegari. The background says that Calegari once conjectured residual finiteness for laminar groups preserving a very full lamination, as an instance of geometric structure restricting group algebra. No separate original citation for Calegari's conjecture is supplied.
\end{itemize}
\end{document}
